% ADASS 2026 proceedings, paper B403 (Birds of a Feather)
\documentclass[11pt,twoside]{article}

% Do NOT use ANY packages other than asp2023.
\usepackage{asp2023}

\aspSuppressVolSlug
\resetcounters

\bibliographystyle{asp2021}

\markboth{Salgado et al.}{Science Platforms at Scale}

\begin{document}

\title{Science Platforms at Scale: Developing Data Services in the Petabyte Era}

\author{Jes\'us~Salgado,$^1$ Christophe~Arviset,$^2$ Javier~Ballester~Arenas,$^2$ Rosie~Bolton,$^1$ Michele~Delli~Veneri,$^1$ Yan~Grange,$^3$ Sandor~Kruk,$^2$ Marcos~L\'opez-Caniego,$^2$ Brian~Major,$^4$ William~O'Mullane,$^5$ and Carmen~Toribio$^6$}
\affil{$^1$SKA Observatory, Jodrell Bank, Macclesfield, UK; \email{jesus.salgado@skao.int}}
\affil{$^2$European Space Astronomy Centre (ESAC), European Space Agency, Villanueva de la Ca\~nada, Madrid, Spain}
\affil{$^3$ASTRON, Netherlands Institute for Radio Astronomy, Dwingeloo, The Netherlands}
\affil{$^4$Canadian Astronomy Data Centre, NRC Herzberg, Victoria, BC, Canada}
\affil{$^5$Vera C. Rubin Observatory, Tucson, AZ, USA}
\affil{$^6$Onsala Space Observatory, Chalmers University of Technology, Onsala, Sweden}

\paperauthor{Jes\'us~Salgado}{jesus.salgado@skao.int}{0000-0002-3680-4364}{SKA Observatory}{}{Macclesfield}{Cheshire}{SK11 9FT}{United Kingdom}
\paperauthor{Christophe~Arviset}{}{}{European Space Agency}{ESAC}{Villanueva de la Canada}{Madrid}{28692}{Spain}
\paperauthor{Javier~Ballester~Arenas}{}{}{European Space Agency}{ESAC}{Villanueva de la Canada}{Madrid}{28692}{Spain}
\paperauthor{Rosie~Bolton}{}{}{SKA Observatory}{}{Macclesfield}{Cheshire}{SK11 9FT}{United Kingdom}
\paperauthor{Michele~Delli~Veneri}{}{}{SKA Observatory}{}{Macclesfield}{Cheshire}{SK11 9FT}{United Kingdom}
\paperauthor{Yan~Grange}{}{}{ASTRON}{}{Dwingeloo}{}{7991 PD}{The Netherlands}
\paperauthor{Sandor~Kruk}{}{}{European Space Agency}{ESAC}{Villanueva de la Canada}{Madrid}{28692}{Spain}
\paperauthor{Marcos~L\'opez-Caniego}{}{}{European Space Agency}{ESAC}{Villanueva de la Canada}{Madrid}{28692}{Spain}
\paperauthor{Brian~Major}{}{}{National Research Council Canada}{Herzberg Astronomy and Astrophysics}{Victoria}{BC}{}{Canada}
\paperauthor{William~O'Mullane}{}{}{Vera C. Rubin Observatory}{}{Tucson}{AZ}{}{USA}
\paperauthor{Carmen~Toribio}{}{}{Chalmers University of Technology}{Onsala Space Observatory}{Onsala}{}{}{Sweden}

%\aindex{Salgado,~J.}
%\aindex{Arviset,~C.}
%\aindex{Ballester~Arenas,~J.}
%\aindex{Bolton,~R.}
%\aindex{Delli~Veneri,~M.}
%\aindex{Grange,~Y.}
%\aindex{Kruk,~S.}
%\aindex{L\'opez-Caniego,~M.}
%\aindex{Major,~B.}
%\aindex{O'Mullane,~W.}
%\aindex{Toribio,~C.}

\begin{abstract}
Science platforms, which bring notebooks, batch processing and visualisation
to the data instead of moving data to users, are now the main way that large
astronomy projects expect their communities to work. This Birds of a Feather
session brought together the teams behind ESA Datalabs and Euclid, the Rubin
Science Platform, the SKA Regional Centre Network (SRCNet), CADC and CANFAR,
NASA Fornax, the LOFAR Long Term Archive and ALMA. Each platform described where
it is today, one design choice it would make again, and its hardest open
problem, with a question for the other platforms. We summarise the answers,
the choices that recur across platforms, and the shared open problems: work
across sites and platforms, efficient use of resources, scaling to petabytes,
identity, and sustaining services through their full lifecycle.
\end{abstract}

%\ssindex{science platform}
%\ssindex{BoFs}
%\ssindex{archives!cloud-based}
%\ssindex{International Virtual Observatory Alliance (IVOA)}

\section{Introduction}

Earlier ADASS discussions \citep{2020ASPC..522..363O, 2020ASPC..527..777T}
established why platforms are needed: data volumes make downloading
impractical, and analysis environments must be close to the archive. Several of
the platforms discussed then are now in production, while new facilities bring
petabytes per year. The focus has therefore moved from \emph{whether} to build a
platform to \emph{how to operate one for years}, and how platforms can work
together. The session opened with a platform round: each platform answered three
prompts (where are you today, what works well, what is the hard problem) on
three slides, followed by themed discussion and an open floor.

\section{Where the Platforms Are Today}

Table~\ref{tab:platforms} summarises the answers to the first prompt. The
platforms span every stage, from a study that has not yet been approved (ALMA)
to services in production for more than a decade (CADC, LOFAR). They differ in
hosting model, from a single data centre to commercial cloud or a federation of
nationally funded centres, but share most building blocks: JupyterLab sessions,
containers on Kubernetes, and IVOA services such as TAP
\citep{2019ivoa.spec.0927D}, ADQL \citep{2023ivoa.spec.1215M} and DataLink
\citep{2023ivoa.spec.1215B}. The platforms are described in more detail
elsewhere: ESA Datalabs \citep{2024sdm..book....1N} and Euclid
\citep{2025AA...697A...1E}, the Rubin Science Platform \citep{2024ASPC..535..227O}
for LSST \citep{2019ApJ...873..111I}, SRCNet \citep{2025ASPC..541...64S},
CADC and CANFAR \citep{2010SPIE.7740E..1IG, 2019ASPC..523..277M}, Fornax
\citep{2024HEAD...2120606J}, the LOFAR archive \citep{2013AA...556A...2V,
2024AC....4800835M} and the ALMA archive \citep{2022Msngr.187...25S}.

\begin{table}[!t]
\caption{Platforms in the round, as reported by their teams in November 2026.}
\label{tab:platforms}
\smallskip
\begin{center}
{\footnotesize
\begin{tabular}{p{2.3cm}p{4.8cm}p{3.9cm}}
\tableline
\noalign{\smallskip}
Platform & Status and users & Data and hosting \\
\noalign{\smallskip}
\tableline
\noalign{\smallskip}
ESA Datalabs & Production; $>$3000 registered users, about 150 active per day; Euclid the main driver & about 5\,PB accessible, growing to 30\,PB; at ESAC \\
Rubin Science Platform & Production; $>$3000 users; 24 deployment environments & LSST; commercial cloud preferred, hybrid planned \\
SRCNet & v0.3 moving into operations; staff and early testers & about 600\,PB per year at full rates; federation of regional centres \\
CADC / CANFAR & Production since 2013; about 900 users, 250 active, 90 projects & 4608 cores, 112 GPUs in the main instance; about 12 instances across SRCNet \\
NASA Fornax & Open beta since January 2026, fully open in January 2027; about 400 users & no own archive; next to NASA data in a commercial cloud region \\
LOFAR LTA & Archive of record since 2010; LOFAR2.0 verification from 2027 & about 50\,PB, over 100\,PB by 2030 \\
ALMA & Platform not yet approved; ALMA-SPACE study by seven European ALMA Regional Centre nodes & 2.7\,PB, growing by about 270\,TB per year \\
\noalign{\smallskip}
\tableline
\end{tabular}
}
\end{center}
\end{table}

\section{Choices the Teams Would Repeat}

Two kinds of choice recurred. The first is \emph{integration with the data}: ESA
Datalabs is built into the ESA science archives, and for Euclid the platform,
the Euclid Archive and ESASky form one environment, so that users who cannot
download pixel data still have discovery, tools and compute in one place. The
second is \emph{standards and reuse}: CANFAR was multi-mission from the start,
built around IVOA standards, and uses the container as the unit of software, so
that the same container runs interactively or in batch. SRCNet reuses Rucio and
FTS \citep{2019CSBS....3...11B} for its data lake and specifies only the
interfaces between independently funded sites. Rubin reported IVOA standards
first, and its Phalanx deployment platform serving 24 environments.

Operational choices were also repeated. Fornax deploys fully automatically, so
the platform can be torn down and redeployed for any upgrade, with no hot fixes;
this kept configuration under control through the beta. LOFAR pointed to
involving the community, which began by necessity when processing LOFAR data
was itself a research problem, and would repeat it in a more planned way, with
clear expectations on credit and sustainability.

\section{Hard Problems and Questions for the Room}

\textbf{Working across sites and platforms.} SRCNet asked how to make jobs run
at independent sites, with the user's identity, quota and data rights, while
feeling like one platform, and who is accountable when they fail. Fornax asked
for a standard that lets a user send a notebook or script from their home
platform to another platform next to a large dataset, run it asynchronously and
retrieve the result. Both point to execution interfaces that do not yet exist
as standards.

\textbf{Using resources efficiently.} CANFAR reported that interactive sessions
can sit idle on reserved resources while infrastructure costs rise. Burstable
jobs increase memory failures, and user feedback from job profiling helps only
partly; vertical scaling and snapshotting idle sessions are untried. The
question was how to use a cluster efficiently without exposing that complexity
to users. Rubin listed operations cost and sustaining the service among its
open issues.

\textbf{Scaling to petabytes.} ESA Datalabs described the work needed to scale
user-friendly workflows to petabyte data, from cutouts over billions of sources
to millions of spectra: hybrid on-premises and cloud infrastructure, NFS to
object storage, FITS to Zarr, Parquet and HATS, and bulk cutout tools. Its
question was where to draw the line between a common platform that can be run
sustainably and the flexibility each mission needs. Rubin is also preparing a
hybrid deployment.

\textbf{Identity.} SRCNet must turn global tokens into local accounts and file
permissions at every site; Rubin uses federated identity through CILogon and
COmanage. Both rely on token-based access that works across institutions.

\textbf{Sustaining services.} LOFAR described the cost of legacy systems after
more than fifteen years of operations: unmanaged service lifecycles make the
infrastructure fragile, and any change breaks some use case. Its question was
how to convince funding agencies that systems need managing and funding
through their complete lifecycle. ALMA showed the other end: interest is high,
but a platform has not yet been approved because of financial and development
constraints, and the ALMA-SPACE study is to set up at least one.

\section{Outlook}

The answers show a common pattern: the hard problems now lie between platforms
and across years, rather than inside a single deployment. Three follow-ups fit
existing community work: execution and job-transfer interfaces between
platforms, taken to the IVOA; shared experience on cost and efficient use of
compute; and arguments for lifecycle funding that platforms can reuse with
their agencies. The organisers thank all platform teams for their answers and
for an open discussion.

\bibliography{B403}

% if we have space left, we might add a conference photograph here. Leave commented for now.
% \bookpartphoto[width=1.0\textwidth]{foobar.eps}{FooBar Photo (Photo: Any Photographer)}

\end{document}
